\section{Wyniki i obliczenia}
\label{sec:wyniki}
\subsection{Czasy pogłosu i kwalifikacja pomieszczenia}
\todo{Wstawić tabelę czasów pogłosu w pasmach oktawowych, z powtórzeniami
oraz opisem konfiguracji pomieszczenia.}
\todo{Zestawić chłonność, pola obu powierzchni pomiarowych, charakterystyki
częstotliwościowe poprawki środowiskowej i wartości A. Porównać metodę
pogłosową z przybliżoną i podać ocenę przydatności pomieszczenia.}

\subsection{Poziomy ciśnienia akustycznego i poprawki}
\todo{Dla każdego położenia źródła zestawić numery i współrzędne punktów,
poziomy źródła z tłem oraz samego tła: A, liniowe i tercjowe.
Podać średnie energetyczne, różnice poziomów i zastosowane poprawki.}

\subsection{Poziom mocy akustycznej i porównanie położeń}
\begin{table}[htbp]
\centering
\caption{Zestawienie końcowe -- bez danych pomiarowych.}
\begin{tabularx}{\linewidth}{@{}L c c c@{}}
\toprule
Położenie źródła & $K_{1A}$ (dB) & $K_{2A}$ (dB) & $L_{WA}$ (dB) \\
\midrule
Środek pomieszczenia & -- & -- & -- \\
Przy ścianie & -- & -- & -- \\
Naroże & -- & -- & -- \\
\bottomrule
\end{tabularx}
\end{table}
\todo{Dołączyć przykładowe obliczenie, wykresy i porównanie wyników.
Brakujących danych nie zastępować zerami.}
